\documentclass[../../../include/open-logic-section]{subfiles}

\begin{document}

\olfileid{sth}{replacement}{absinf}
\olsection{પ્રતિસ્થાપન અને ``નિરપેક્ષ અનંતતા''}

હવે \limofsize{}ને બાજુએ મૂકીને પ્રતિસ્થાપનને વાજબી
ઠરાવવાના બીજા પ્રકારના (આંતરિક) પ્રયાસો જોઈશું. તેમાં
ગણો તબક્કાઓમાં રચાય છે એ વિચારને ગંભીરતાથી લેવાય છે.

પહેલી વાર પુનરાવર્તી પ્રક્રિયાની રૂપરેખા આપતાં, દરેક
તબક્કે શું થાય છે તે સમજાવતા કેટલાક સિદ્ધાંતો આપ્યા
હતા: \stageshier, \stagesord અને \stagesacc. પછી
તબક્કાઓની સંખ્યા વિશે કંઈક જણાવતા સિદ્ધાંતો ઉમેર્યા:
\stagessucc{} કહે છે કે ગણરચનાની પ્રક્રિયા કદી પૂરી
થતી નથી; અને \stagesinf{} કહે છે કે પ્રક્રિયા એક
અનંત ક્રમાંકના તબક્કે પહોંચે છે.

આ છેલ્લાં બે સિદ્ધાંતો નીચેના જેવા કોઈ વધુ વ્યાપક
સિદ્ધાંતમાંથી નીકળે એવું સૂચવવું વાજબી છે:
\begin{enumerate}
	\item[] \stagesinex. તબક્કાઓની સંખ્યા નિરપેક્ષ રીતે
	અનંત છે; પદાનુક્રમ શક્ય હોય એટલો ઊંચો છે.
\end{enumerate}
દેખીતી રીતે, આ અનૌપચારિક સિદ્ધાંત છે. પરંતુ ગણની
સંચયી-પુનરાવર્તી કલ્પનામાંથી તે તરત જ \emph{અનિવાર્યપણે
નીકળતો} ન હોય તો પણ તેની સાથે \emph{સુસંગત} જરૂર
લાગે છે. ઓછામાં ઓછું, \limofsize{}થી વિપરીત, તે
ગણો તબક્કે તબક્કે રચાય છે એ વિચાર જાળવે છે.

હવે આશા એ છે કે \stagesinex{}ના આધારે પ્રતિસ્થાપનને
વાજબી ઠરાવી શકાય. તે કેવી રીતે શક્ય હોય તે જોઈએ.

\olref[ordinals][idea]{sec}માં જોયું કે (સહજ અર્થમાં)
$\omega+\omega$ સાથે ક્રમ-એકરૂપ હોવો જોઈએ એવો સુક્રમ
રચવો સરળ છે. બીજી રીતે કહીએ તો, જેનો ક્રમપ્રકાર
(સહજ અર્થમાં) $\omega+\omega$ હોય એવી તબક્કાવાર
પુનરાવર્તી પ્રક્રિયા સહેલાઈથી કલ્પી શકાય છે. તેથી
\stagesinex{} સ્વીકાર્યું હોય, તો પદાનુક્રમમાં ઓછામાં
ઓછો $\omega+\omega$મો તબક્કો, એટલે
$V_{\omega+\omega}$, છે એવું પણ નિશ્ચિતપણે સ્વીકારવું
જોઈએ; પદાનુક્રમ એટલા સુધી \emph{ચાલી શકે} જ છે.

આ વિચારને સામાન્ય બનાવીએ: કોઈ પણ સુક્રમ માટે
પુનરાવર્તી પદાનુક્રમ રચવાની પ્રક્રિયા ઓછામાં ઓછી તે
સુક્રમ જેટલી આગળ જવી જોઈએ. અને આની ખાતરી માટે
\olref[ordinals][ordtype]{thmOrdinalRepresentation}ને
\emph{સ્વયંસિદ્ધિ} તરીકે લેવું જ પૂરતું છે. તે જણાવશે
કે દરેક સુક્રમ કોઈ વૉન નોયમન ક્રમવાચક સંખ્યા સાથે
ક્રમ-એકરૂપ છે. દરેક વૉન નોયમન ક્રમવાચક સંખ્યા તેના
પોતાના સ્તરાંક જેટલી હોવાથી,
\olref[ordinals][ordtype]{thmOrdinalRepresentation}
પછી બતાવશે કે ગણસિદ્ધાંતમાં જ્યારે પણ કોઈ સુક્રમનું
વર્ણન કરી શકીએ, ત્યારે ગણરચનાની પુનરાવર્તી પ્રક્રિયા
તે સુક્રમથી આગળ જવી પડે છે.

આ વિચાર ચોક્કસપણે \stagesinex{}ના ફલિતાર્થ જેવો લાગે
છે. દુર્ભાગ્યે, આ વિચારમાંથી પ્રતિસ્થાપન મેળવવાનો હેતુ
હોય, તો એક સરળ તકનિકી અવરોધ છે: પ્રતિસ્થાપન
\olref[ordinals][ordtype]{thmOrdinalRepresentation}
કરતાં કડક રીતે વધુ શક્તિશાળી છે. (આ વાત
\citet[\S13.2]{Potter2004}એ નોંધેલી છે; આપણે તેને
\olref[sth][replacement][finiteaxiomatizability]{sec}માં
સાબિત કરીશું.)

આથી \stagesinex{}ને પ્રતિસ્થાપન આપે એવા અર્થમાં
સમજવું હોય, તો તે \emph{માત્ર} દરેક સુક્રમથી
પદાનુક્રમ આગળ જાય છે એટલું કહી શકે નહીં. તેનાથી
વધુ શક્તિશાળી દાવો જોઈએ. આ માટે \cite{Shoenfield:AST}એ
વિચારનું બહુ સ્વાભાવિક મજબૂત સ્વરૂપ સૂચવ્યું: પદાનુક્રમ
કોઈ પણ ગણ સાથે \emph{સહઅંતિમ} નથી.\footnote{ગોડેલે
પણ કદાચ આવો જ વિચાર સૂચવ્યો હતો; જુઓ
\citet[p.~223]{Potter2004}. ગોડેલ અને
\citeauthor{Shoenfield:AST}ની ચર્ચા માટે
\citet[90--5]{Incurvati2020} જુઓ.} થોડું વિગતે:
જો $\tau$ ગણોને પદાનુક્રમના તબક્કાઓમાં મોકલતું
પ્રતિચિત્રણ હોય, તો કોઈ પણ ગણ $A$ની $\tau$ હેઠળની
પ્રતિમા આખા પદાનુક્રમને પૂરી કરતી નથી. બીજી રીતે
કહીએ તો (પ્રરૂપાત્મક રીતે):
\begin{enumerate}
	\item[] \stagescofin. જો $A$ ગણ હોય અને દરેક
	$x \in A$ માટે $\tau(x)$ તબક્કો હોય, તો દરેક
	$\tau(x)$ પછી આવતો એક તબક્કો છે, જ્યાં $x \in A$.
\end{enumerate}
યોગ્ય રીતે ઔપચારિક બનાવેલું \stagescofin{}નું સ્વરૂપ
$\ZF$ સાબિત કરે છે તે સ્પષ્ટ છે. ઊલટી તરફ,
\stagescofin{} પ્રતિસ્થાપનને વાજબી ઠરાવે છે એવી
અનૌપચારિક દલીલ આપી શકાય.\footnote{\stagescofin{}ના
કોઈ ઔપચારિક સ્વરૂપથી પ્રતિસ્થાપન સાબિત કરવું વધુ મુશ્કેલ
થશે, કારણ કે એકલું $\Z$ તબક્કાઓ વ્યાખ્યાયિત કરવા
પૂરતું શક્તિશાળી નથી. તેથી \stagescofin{}નું
ઔપચારિક સ્વરૂપ કેવી રીતે આપવું તે સ્પષ્ટ નથી. તેમ છતાં,
\olref[sth][replacement][extrinsic]{sec}માં ચર્ચ્યા મુજબ
$\LT$ના કોઈ વિસ્તરણમાં કામ કરવું એક વિકલ્પ છે.}
ખરેખર, ધારો કે $(\forall x \in A)\lexists![y][\phi(x,y)]$.
પછી દરેક $x \in A$ માટે $\sigma(x)$ એ
$\phi(x,y)$ને સંતોષતું $y$ લો, અને $\tau(x)$ એ
$\sigma(x)$ પહેલી વાર રચાય તે તબક્કો લો.
\stagescofin{} મુજબ એવો તબક્કો $V$ છે કે
$(\forall x \in A)\tau(x)\in V$. હવે દરેક
$\tau(x) \in V$ અને $\sigma(x) \subseteq \tau(x)$
હોવાથી, પૃથક્કરણ વડે
$\Setabs{y
\in V}{(\exists x \in A)\sigma(x) = y} = \Setabs{y}{(\exists x \in
A)\phi(x,y)}$ મેળવી શકાય છે.

\begin{prob}
	$\ZF$ની અંદર \stagescofin{}ને ઔપચારિક સ્વરૂપ આપો.
\end{prob}

આમ \stagescofin{} પ્રતિસ્થાપનને વાજબી ઠરાવે છે.
અને \stagesinex{}થી \stagescofin{} વાજબી ઠરે તે
ઓછામાં ઓછું સંભવિત લાગે છે. કારણ કે ધારો કે
\stagescofin{} નિષ્ફળ જાય. ત્યારે પદાનુક્રમ કોઈ ગણ~$A$
સાથે સહઅંતિમ છે: એટલે કે એવું પ્રતિચિત્રણ $\tau$ છે
કે કોઈ પણ તબક્કા~$S$ માટે કોઈ $x \in A$ એવો મળે
કે $S \in \tau(x)$. આ સ્થિતિમાં પદાનુક્રમની કથિત
``નિરપેક્ષ અનંતતા''ને પકડી લેવાની એક રીત મળે છે:
તે $A$ પર લાગુ કરેલા $\tau$ની વ્યાપ્તિથી
\emph{પૂરી થાય છે}. અને તેથી પદાનુક્રમ ``નિરપેક્ષ
રીતે અનંત'' છે એ વિચાર નબળો પડે છે. તેનો પ્રતિપક્ષ
લઈએ તો: \stagesinex{}માંથી \stagescofin{} નીકળે,
અને તે પછી પ્રતિસ્થાપનને વાજબી ઠરાવે છે.

આ પ્રતિસ્થાપન માટે આંતરિક ઔચિત્ય આપવાનો ખરેખર
આશાસ્પદ પ્રયાસ છે. પરંતુ છેવટે તે સફળ થાય છે કે
નહીં, એ નક્કી કરવાનું તમારા પર છોડીશું.

\end{document}
